\subsection{The sharp law at every fixed first-layer rank}
\label{sec:fixed-rank-batching}

Rectangular multiplication removes the rank-three restriction in the
preceding preprocessing argument. The analytic construction already works
at every fixed rank. Its Taylor tables have a number of columns that is
polynomial in the logarithm of the width. A sufficiently unbalanced exact
matrix multiplication computes all those columns within the original
preprocessing budget.

\begin{theorem}[Critical sampling at every fixed first-layer rank]
\label{thm:fixed-rank-critical}
Fix an integer $r\ge1$. Consider a zero-bias width-$n$, depth-$D\ge1$
tanh network whose absolute row sums are at most one and whose first
weight matrix has rank at most $r$. Its weights are dyadic with at most
$b$ bits. Put $B=16+b+\lceil\log_2(nD+2)\rceil$.
One uniform algorithm, with constants depending on $r$, samples the
exact output using
\[
 O_r(\sqrt{D+1})\quad\hbox{expected input probes},\qquad
 O_r\!\left(\sqrt{D+1}\,B^{\max\{4,r+2\}}\right)
 \quad\hbox{expected bit operations}.
\]
At $D=\lceil\log_2n\rceil$ and $b=O(\log n)$, its weight-only
preprocessing and storage are $O_r(Dn^2\log^6n)$.
The worst-case query order is $\Theta_r(\sqrt D)$ when
$n\ge D+1$ and $b\ge\lceil\log_2(D+1)\rceil$.
\end{theorem}

The rank is fixed before taking the asymptotic bound. Both the constants
and the degree of the online precision polynomial depend on it. The
theorem supplies no uniform polynomial bound when $r$ grows with $n$.
In particular, rank four is not a hardness threshold.

\begin{lemma}[Exact batching of a thin coefficient matrix]
\label{lem:fixed-rank-rectangular-batching}
For every fixed integer $r\ge1$, an $n$ by $n$ integer matrix can be
multiplied by an $n$ by $C$ integer matrix, $C\ge1$, using
\[
 O_r\!\left(\log(C+2)
 \left[n^2C^{1/(2r)}+C^{8+1/(2r)}\right]\right)
 \tag{*}\label{eq:fixed-rank-rectangular-work}
\]
scalar arithmetic operations. If the input entries have at most $m$
bits, the computation can use only integers with
$O_r(m+\log n+\log(C+2))$ bits. It is a uniform exact algorithm;
no oracle for real arithmetic or for matrix multiplication constants
is required.
\end{lemma}

\paragraph{Proof idea.}
Choose one rational bilinear formula with a square side $q$ and a
short side $w\ge q^{1/4}$. Its tensor powers supply a block large
enough to hold all $C$ coefficient columns. Tiling the large square
matrix with these blocks yields the first term in
\eqref{eq:fixed-rank-rectangular-work}; padding yields the second.

\begin{proof}
Set $\epsilon=1/(8r)$. Le Gall's rectangular multiplication bound
gives exponent two at every aspect exponent below $0.3029805$
\cite{legall2012rect}. We use the strictly smaller exponent $0.30$.
The theorem and the definition of the exponent by exact tensor rank
apply over the rational field. Permuting the three tensor indices
gives multiplication of a $q$ by $q$ matrix by a $q$ by $w$ matrix.
Consequently there are integers $q\ge2$, $2\le w\le q$, and a finite
exact rational bilinear formula with $R$ products such that
\[
 w^4\ge q,\qquad q^2\le R\le q^{2+\epsilon}.
 \tag{**}\label{eq:fixed-rank-base-parameters}
\]
To see that the inequalities can hold simultaneously, apply the
exponent bound with excess $\epsilon/2$, then take a sufficiently
large base size and $w=\lfloor q^{0.30}\rfloor$.
The constant in the rank bound is absorbed by $q^{\epsilon/2}$.
The tensor flattening along the $q^2$ entries of the first input
has rank $q^2$, which gives $R\ge q^2$.

There is also a direct uniform way to choose this finite formula.
Enumerate integer triples and rational coefficient arrays by increasing
encoding length. For each candidate, check
$w^4\ge q$ and $R^{8r}\le q^{16r+1}$, and verify every coefficient
in its bilinear tensor identity by rational arithmetic. This search
terminates by the preceding existence theorem. It is performed once
for a given $r$; its finite time and stored certificate are constants
in $O_r$. Thus this use of the exponent theorem does not hide
nonuniform advice or a real coefficient test.

Fix the resulting formula. A tensor power of depth $j$ multiplies
a $u$ by $u$ matrix by a $u$ by $c$ matrix, where
$u=q^j$ and $c=w^j$. Its operation count $T_j$ satisfies
\[
 T_j\le R T_{j-1}+O_r(q^{2j}+q^jw^j)
       \le R T_{j-1}+O_r(q^{2j}).
\]
Since $q^2\le R\le q^{2+\epsilon}$, this implies
\[
 T_j=O_r\bigl((j+1)u^{2+\epsilon}\bigr).
\]
The factor $j+1$ includes the possible equality $R=q^2$.

Choose the least $j\ge0$ such that $w^j\ge C$.
Then $c\le wC$, $j+1=O_r(\log(C+2))$, and
\[
 u\le q C^{\log q/\log w}\le qC^4.
\]
Pad the row dimension to a multiple of $u$ and the column dimension
to $c$. With $h=\lceil n/u\rceil$, there are $h^2$ block
products. Their output sums and the creation of zero padding fit the
same operation count. Using $h\le n/u+1$ gives
\[
 O_r\bigl((j+1)[n^2u^\epsilon+u^{2+\epsilon}]\bigr)
 =
 O_r\!\left(\log(C+2)
 [n^2C^{4\epsilon}+C^{8+4\epsilon}]\right).
\]
This argument also covers $C>n$. Substitution of
$4\epsilon=1/(2r)$ proves
\eqref{eq:fixed-rank-rectangular-work}.

It remains to justify the bit lengths. Let $\Delta$ be a positive
common denominator of the finitely many rational coefficients in the
chosen base formula. Every input transformation through $j$ levels
has denominator dividing $\Delta^j$. Leaf products therefore have
denominator dividing $\Delta^{2j}$, and the $j$ output
recombination levels raise this exponent by at most $j$.
All intermediate quantities can be represented over predetermined
powers of $\Delta$ of exponent at most $3j$.
No reduction by a greatest common divisor is necessary.

The base formula contains only a fixed number of fixed coefficients.
Their absolute linear-form sums give a fixed bound on magnitude
growth at one level. After input transformation, multiplication,
and output recombination, the logarithm of the magnitude has grown
by $O_r(j)$, in addition to the input bit lengths. Summation over
the $h$ inner blocks adds at most $O(\log n)$ bits. Hence integer
numerators, fixed-denominator rescalings, and completed output
entries all have $O_r(m+j+\log n)$ bits. Applying schoolbook
arithmetic to these integers proves the stated finite-bit
implementation. The identities hold exactly before any later
rounding of the completed matrix product.
\end{proof}

\paragraph{Proof idea.}
The proof of Theorem~\ref{thm:rank-three-critical} builds its factory
and analytic table for every fixed number of scalar source streams;
only its matrix multiplication estimate uses rank three. Keep the
factory, its amplitude gate, and its precision stopping rule. Apply
the batching lemma at each dense linear layer, then count the
remaining coefficient arithmetic directly. For fixed rank, the extra
powers of $\log n$ in the latter work are absorbed by its smaller
factor of $n$.

\begin{proof}[Proof of Theorem~\ref{thm:fixed-rank-critical}]
If the first matrix has rank zero, the output has mean zero and a
fair bit suffices. Otherwise use its actual rank, padding the scalar
source description to $r$ if desired. The construction before the
matrix multiplication step in the proof of
Theorem~\ref{thm:rank-three-critical} applies to every fixed $r$.
It gives an amplitude gate of probability $O(D^{-1/2})$, a
conditional factory using $O_r(D+1)$ expected probes, and a Taylor
oracle whose conditional expected bit work is
$O_r((D+1)B^{\max\{4,r+2\}})$.
That proof includes the unconditional probability bounds for all
arbitrary precision refinements. These facts yield the claimed
online bounds. Depth one uses the direct scalar factory.

The table of that proof has $J=O_r(D^{r/2})$ grid centers, Taylor
degree $K=O_r(B)$, $N_K=\binom{K+r}{r}=O_r(K^r)$ coefficients per
center, $C=JN_K=O_r(D^{r/2}K^r)$ coefficient columns, and working
precision $M=O_r(B+D)$. Increasing $M$ by
$O_r(b+\log n+\log(C+2))$ preserves this order.

At each linear layer, clear denominators as in that proof and apply
Lemma~\ref{lem:fixed-rank-rectangular-batching} to the resulting
integer matrices. All intermediate integers have $O_r(M)$ bits,
so each scalar arithmetic operation costs $O_r(M^2)$ bit
operations. The product is exact before its completed entries are
rounded, so the analytic error induction applies without change.
The basis computation and the nonlinear recurrences have the costs
given in that proof. A general-depth preprocessing bound is thus
\[
 O_r\!\left(
 n^2B^3+
 D\left[
 \log(C+2)\bigl(n^2C^{1/(2r)}+C^{8+1/(2r)}\bigr)
 +nD^{r/2}K^{2r}+nD^{r/2}M
 \right]M^2\right).
 \tag{***}\label{eq:fixed-rank-general-preprocessing}
\]
This bound includes the first affine layer, all working storage,
and the stored Taylor table.

Let $v=\log_2 n\ge1$. When $D=\lceil\log_2n\rceil$ and
$b=O(\log n)$, the chosen parameters satisfy
$D,K,M=O_r(v)$ and $C=O_r(v^{3r/2})$.
The dense term in
\eqref{eq:fixed-rank-general-preprocessing} is at most
\[
 O_r\!\left(n^2v^{3+(3r/2)/(2r)+1}\right)
 =O_r(n^2v^{19/4}),
\]
where the final factor $v$ bounds $\log(C+2)$.
The nonlinear term is $O_r(nv^{3+5r/2})$, the constant
exponentials cost $O_r(nv^{4+r/2})$, and the padding term is
$O_r(v^{12r+19/4})$. For each fixed $r$, each of the last three
terms is $O_r(n^2v^7)$, because a fixed power of $\log n$ is
bounded by any fixed positive power of $n$. The basis term also
fits this bound. Consequently preprocessing and storage are
$O_r(n^2\log^7n)=O_r(Dn^2\log^6n)$.

Finally, the rank-one mean-chain family lies in every class with
rank bound $r\ge1$. Its lower bound from
Theorem~\ref{thm:finitecritical} is
$\Omega(\sqrt D)$ under the displayed width and dyadic-grid
conditions. This proves the worst-case claim.
\end{proof}

The rare full-input branch in
Corollary~\ref{cor:rank-three-finite-input} uses the same Taylor
oracle and hence also extends to every fixed $r$. Under
\eqref{eq:fixed-rank-general-preprocessing}, its expected probe
bound is $O_r(\min\{\sqrt D,n/\sqrt D\})$, with the same
finite-input lower bound and with expected bit work
$O_r((1+\min\{\sqrt D,n/\sqrt D\})B^{\max\{4,r+2\}})$.
The unrestricted-rank critical exponent remains open.
